% Shared preamble for the standalone circuit figures in this directory.
%
% qcircuit runs on xy-pic, which costs ~26ms per diagram on every compile of the
% thesis; drawn once here and included as PDFs, they cost the thesis nothing.
% The one circuit still inline in the main text (Figure 4.1) keeps qcircuit in
% the thesis preamble -- it sits among text that reflows, and is edited more
% often than these.
%
% Eleven drawings share this preamble: the five card circuits of Appendix D's
% Figures D.1-D.5 (card_dec_*.tex), the five reference drawings they were built
% from (circuit_dec_*.tex, read by code/_build_povm_cards.py and
% code/_build_recipe_figure.py rather than included), and Figure 1.1's
% mini-circuit.  The macros they share are declared here once and copied from
% bsc-thesis.tex; only the handful they use is here, and it must keep agreeing.
%
% 11pt matches the thesis document class, so a circuit built here is the same
% size it was inline; the figures are therefore included at natural size, with
% no width= to rescale them.
%
% The border is {left bottom right top} and every number in it was measured,
% not guessed. xy-pic under-declares its bounding box: the \lstick labels
% (|0>, \rho) hang 13.5pt off the left of it and the outer wires another 0.4pt
% -- half a rule width -- off each remaining side. At border=0pt the labels are
% silently CLIPPED AWAY. Hence 14pt and 0.5pt.
%
% Left and right are equal on purpose. Inline, the box being centred was the
% one xy declared, i.e. the wires without their labels; padding both sides
% equally keeps that same box centred, so the circuits land exactly where they
% used to on the page rather than shifting ~7pt.
%
% Re-measure if a circuit ever changes shape: build at a large known border,
% ask `gs -sDEVICE=bbox` for the true ink box, and the difference per side is
% what the border has to cover.
\documentclass[11pt, border={14pt 0.5pt 14pt 0.5pt}]{standalone}

\usepackage{mathtools}  % loads amsmath (smallmatrix, \text)
\usepackage{amssymb}
\usepackage{bbm}
\usepackage{qcircuit}

\newcommand{\ket}[1]{|#1\rangle}
\newcommand{\Id}{\mathbbm{1}}
\newcommand{\phig}{\tau}   % golden ratio, Conway & Smith
